\chapter{LITERATURE REVIEW}

\section{Development of a new signal processing algorithm based on independent component analysis for single channel ECG data \cite{paper_1}}
Lee, J., Lee, K. J., and Yoo, S. K. (XXXX) proposed a novel signal processing algorithm based on Independent Component Analysis (ICA) for single-channel ECG data enhancement and QRS complex detection.

The main objective was to overcome the limitation that ICA typically requires multi-channel signals by developing a method to apply ICA to single-channel ECG recordings. The key contribution was demonstrating that signal enhancement could be achieved even with single-channel data by constructing artificial multi-channel signals through time-delay embedding.

The technical approach involved generating 1-20 sample point delayed versions of the original single-channel ECG to create a 20-channel constructed signal. The FastICA (fixed-point algorithm) was then applied to extract independent source signals. For QRS detection validation, Hilbert transform was used to compute analytic signals, from which instantaneous phase and instantaneous frequency were derived. Wavelet transform was subsequently applied to the instantaneous frequency signal, using a level 3 approximation for final QRS detection.

The data used in this study was obtained from the MIT/BIH arrhythmia database, a standard benchmark for ECG algorithm evaluation.

Performance evaluation focused on QRS complex detection efficacy using the enhanced signals. The study specifically examined records that could not be properly filtered using conventional techniques (records 108 and 117), demonstrating the algorithm's ability to handle challenging noise conditions including motion artifacts and power line interference.

Key findings revealed that the delay parameter significantly affects performance: too many delayed sample points (large channel count) caused signal smoothing, while too few reduced enhancement effectiveness. The optimal recommendation was 10-15 point delays for signals sampled at 360 Hz. The proposed algorithm successfully removed noises that conventional filtering could not handle, and QRS complex detection was perfectly accomplished using the instantaneous frequency signal of the enhanced output. The enhanced ECG signals were found profitable for clinical data reading.

The authors concluded that the ICA-based single-channel signal processing algorithm showed benefits in signal reading and QRS detection. Future work was identified as algorithm optimization and simplification for practical implementation.

This paper is highly relevant to our capstone project as it addresses the fundamental challenge of processing single-channel cardiac signals—similar to our smartphone-based acquisition where we extract cardiac information from a single PPG or SCG channel. The time-delay embedding technique for enabling ICA on single-channel data could potentially be applied to enhance our mobile-acquired signals before peak detection. Furthermore, the paper's approach to handling motion artifacts and noise that conventional filters cannot address directly parallels our challenge of processing noisy smartphone sensor data. The Hilbert transform and instantaneous frequency analysis methods described could inform our Stage 2 DSP pipeline for robust systolic peak detection.


\newpage
\section{Cardiac Conduction Model for Generating 12 Lead ECG Signals With Realistic Heart Rate Dynamics \cite{paper_2}}
Quiroz-Juárez, M. A., Jiménez-Ramírez, O., Vázquez-Medina, R., Ryzhii, E., Ryzhii, M., and Aragón, J. L. (2019) proposed an extended heterogeneous oscillator model of the cardiac conduction system for generating realistic 12-lead ECG waveforms with heart rate variability (HRV). The main objective was to develop a biophysically grounded model capable of producing all standard 12 ECG leads while incorporating realistic heart rate dynamics, including Mayer waves, respiratory sinus arrhythmia (RSA), additive noise, and baseline wander, and to fit the model to individual subject data. The key contribution was a unified framework that combines modified Van der Pol oscillators for the natural pacemakers (SA, AV, and His-Purkinje nodes) with modified FitzHugh-Nagumo oscillators for atrial and ventricular muscles, enabling simulation of both normal and pathological ECG patterns (e.g., Wellens syndrome).

The technical approach used coupled ordinary differential equations: the pacemakers were modeled as modified Van der Pol oscillators with time-delay velocity coupling, while the atrial and ventricular depolarization/repolarization processes were described by modified FitzHugh-Nagumo equations. An artificial RR-tachogram was generated with a bimodal power spectrum (Mayer wave at 0.1 Hz and RSA at 0.25 Hz) and incorporated into the model by dynamically varying the intrinsic frequency parameter f1f1​ of the SA node. The 12-lead ECG was then computed as a weighted linear combination of the atrial and ventricular responses, with weighting coefficients obtained via a perceptron-like supervised learning algorithm that fitted the model output to real ECG signals. Baseline wander coupled to respiratory frequency and additive Gaussian noise were also included.

The data used for fitting and validation included real 12-lead ECG recordings from the PhysioNet database (13 subjects: 9 men, 4 women, aged 26–69) and clinical ECG examples from published libraries (e.g., Wellens syndrome Type 1 ischemia). Performance evaluation was based on the mean absolute error between the model-generated and real ECG waveforms, with cumulative average error across subjects demonstrating asymptotic convergence. Qualitative comparisons of the 12-lead waveforms and pathological patterns (ST-segment elevation, T-wave inversion) were also presented.

Key findings showed that the model accurately reproduced realistic ECG characteristics, including correct RR and QT interval relationships, HRV spectral peaks (LF and HF bands), and lead-specific morphologies consistent with Einthoven’s triangle. The perceptron-based fitting procedure converged within a finite number of iterations and yielded low mean absolute errors. The model successfully simulated pathological conditions such as Wellens syndrome, demonstrating its potential for medical education and ECG signal processing system testing. The authors concluded that the proposed model provides a useful tool for generating realistic 12-lead ECG signals with heart rate dynamics and biophysical properties, and suggested future work on developing a wide-range ECG simulator and incorporating more complex pathological phenomena.

This paper is directly relevant to our Stage 4 (Member D – Dynamics) digital twin simulation. The heterogeneous oscillator architecture—using modified Van der Pol equations for pacemakers and FitzHugh-Nagumo equations for muscle tissue—closely mirrors the coupled nonlinear cardiac oscillator model we plan to implement for simulating rhythm breakdown under autonomic stress. The paper’s method for incorporating HRV statistics (Mayer waves and RSA) into the oscillator frequency parameter f1f1​ offers a concrete strategy for mapping our extracted PRV features (LF/HF ratio, SDNN) into model parameters (damping μμ, coupling c12c12​). Furthermore, the perceptron-based parameter fitting approach could be adapted for personalizing the digital twin to individual users based on their measured pulse rate variability. Finally, the inclusion of baseline wander and noise models provides insight into realistic signal degradation, which is relevant to our Stage 2 DSP pipeline for robust peak detection in noisy smartphone-acquired signals.
